%% 致谢

\begin{acknowledgements}
在论文即将完成之际，回顾本次毕业设计从选题确定、系统实现、实验分析到论文撰写的全过程，我深切感受到本课题的顺利完成离不开老师和同学们在学习、科研与实践过程中给予的指导、帮助和支持。正是在他们的关心与鼓励下，我才能够在遇到困难时不断调整思路、完善方案，并最终完成本次毕业设计。在此，我谨向所有关心和帮助过我的老师和同学表示诚挚的感谢。

首先，衷心感谢我的指导教师戴玉超老师。在课题开展初期，老师帮助我明确研究方向，使我能够围绕深空场景的三维重建与新视角合成展开系统设计；在方法实现和实验分析过程中，老师多次指出方案中存在的问题，并从研究思路、论文结构、实验设计和结果表达等方面给予了细致指导。老师严谨的治学态度和认真负责的工作作风，使我在完成毕业设计的同时，也更加理解了科研工作中逻辑严密、论证充分和实事求是的重要性。

同时，诚挚感谢学院的各位老师。在大学的学习时光里，老师们凭借丰富的教学经验和扎实的专业知识，认真讲授每一门课程，耐心解答学习过程中遇到的问题，使我逐步建立起较为系统的专业知识体系。从基础理论的学习到专业能力的培养，老师们的悉心教导为我后续开展毕业设计和论文写作奠定了坚实基础。课堂上的讲解、课后的指导以及严谨负责的教学态度，都让我受益匪浅。

感谢同学们在课题推进过程中给予的帮助和建议。在代码调试、实验运行和论文修改过程中，大家的交流与讨论使我少走了许多弯路，也让我在遇到问题时能够从不同角度思考解决方案。与同学们共同学习、相互鼓励的过程，不仅帮助我顺利推进了毕业设计，也让我更加认识到交流合作在学习和科研中的重要作用。

最后，感谢评审专家和答辩委员会老师在百忙之中审阅本文并提出宝贵意见。毕业设计的完成不是学习的终点，而是新的开始。今后我将带着本次课题中获得的经验与反思，继续保持认真严谨的态度，在后续学习和工作中不断提升自己的专业能力。
\end{acknowledgements}
